% ID: SEL-D0055-OB-E01
% ESTADO: cotejada. Segunda pasada visual y borrador compilado revisados.
% Fuente: Selectividad/sel_2011_matematicas_aplicadas/122A-EXAMEN 6.pdf
% SHA-256: aebcbfddd4d5db5f5a10b4662e9e12b2ee362bbf574b751bccd6d78594766ad3
% Página: 2; opción B; ejercicio original 1.

\begin{enumerate}[label=\alph*)]
\item \textbf{(1.5 puntos)} Dadas las matrices
\[
M=\begin{pmatrix}0&3&-1\\1&0&-2\end{pmatrix}\quad\text{y}\quad
N^t=\begin{pmatrix}2&3&-1\\-1&1&0\end{pmatrix},
\]
razone cuáles de las siguientes operaciones tienen sentido y efectúe las que puedan realizarse:
\[
M+N^t,\quad M^t\cdot N,\quad M\cdot N.
\]
\item \textbf{(1 punto)} Un industrial cafetero produce dos tipos de café, natural y descafeinado, en tres modalidades cada uno, A, B y C. Se han anotado en la matriz $P$ los pesos, en kg, del café que el industrial produce de cada una de las modalidades de cada tipo, y en la matriz $Q$ los precios a los que vende el kg de cada producto final:
\[
P:\quad\bordermatrix{&A&B&C\cr\text{natural}&550&400&240\cr\text{descafein.}&260&200&100\cr}
\]
\[
Q:\quad\bordermatrix{&A&B&C\cr\text{natural}&2.20&2.75&2.50\cr\text{descafein.}&3.20&3.90&3.60\cr}
\]
Efectúe el producto $P\cdot Q^t$ y explique el significado económico de cada uno de los elementos de la diagonal principal de la matriz resultante.
\end{enumerate}
